% cap02/2.4_dds.tex
\section{Evaluación y Descarte Arquitectónico del Estándar DDS (\textit{Data Distribution Service})}

El rigor metodológico exige documentar no solo las tecnologías adoptadas, sino
también el descarte justificado de aquellas que resultan inviables para el
ecosistema objetivo. En este proyecto, orientado a las ONG, el estándar
\textit{Data Distribution Service} (DDS) de la \textit{Object Management Group}
(OMG) constituye un caso ilustrativo de una tecnología de altísima precisión cuyo
traslado a un entorno civil, móvil y de conectividad intermitente resulta
contraproducente.

\subsection{Conceptualización y Origen del Estándar DDS}

DDS fue concebido en el sector de defensa y aeronáutica para gobernar el intercambio
de telemetría en tiempo real. Su diseño se basa en el patrón de mensajería
desacoplada \textit{Data-Centric Publish-Subscribe} (DCPS), que eleva el dato a
entidad de primera clase dentro de la topología del sistema. Este esquema se apoya
en el protocolo \textit{Real-Time Publish-Subscribe} (RTPS), optimizado para operar
sobre datagramas ligeros (principalmente UDP y, en ciertos despliegues, TCP). Su
objetivo es que numerosos publicadores y suscriptores se descubran de forma
descentralizada, sin depender de un intermediario o \textit{broker} central.

\subsection{El Choque Paradigmático: Entorno Militar frente al Ecosistema Civil SaaS}

Al evaluar la factibilidad de trasladar esta topología a la plataforma
\textit{Democra} surgieron dos obstáculos decisivos. Primero, los cortafuegos
(\textit{firewalls}) públicos y corporativos: DDS depende del ruteo por múltiples
puertos dinámicos para su descubrimiento entre pares, un tráfico que las redes
públicas ---la conexión móvil de un voluntario o el \textit{WiFi} restringido de una
sede rural--- suelen bloquear, dejando la aplicación inoperante. Segundo, el modelo
descentralizado sin \textit{broker} contradice una necesidad esencial de
\textit{Democra}: la auditoría centralizada de cada transacción. El tercer sector no
requiere que los dispositivos se comuniquen entre sí de forma directa, sino que cada
cliente se conecte a un servidor central en la nube (\textit{Node.js}) que procese,
persista y audite la información en \textit{PostgreSQL} bajo un modelo
\textit{multi-tenant}.

\subsection{El Triunfo Arquitectónico de los WebSockets sobre DDS}

La decisión de descartar DDS condujo a adoptar el protocolo estándar y nativo de la
web, \textit{WebSockets} (RFC 6455). A diferencia del tráfico UDP de DDS, los
\textit{WebSockets} operan sobre el puerto seguro y universal HTTPS/TLS (443), lo que
les permite atravesar sin dificultad los cortafuegos de las redes públicas y
rurales. Con ello se logra una comunicación bidireccional en tiempo real que permite
a la interfaz \textit{React.js} de \textit{Democra} actualizar los inventarios de las
distintas organizaciones sin recurrir a protocolos de origen militar ajenos al
ecosistema web.
